% Part: first-order-logic
% Chapter: introduction
% Section: substitution

\documentclass[../../../include/open-logic-section]{subfiles}

\begin{document}

\olfileid{fol}{int}{sub}

\olsection{आदेशन}

सर्व गोष्टींचा परस्परसंबंध कसा आहे आणि विन्यासमीमांसा व
चिन्हार्थमीमांसा यांचा विकास पुढील अधिक प्रगत अभ्यासाचा पाया कसा घालतो,
हे दाखवण्यासाठी एका उदाहरणाची चर्चा करू. आपल्या !!{derivation} पद्धतींनी
आपल्याला $\lforall[\Obj v_0][\Atom{\Obj P}{\Obj v_0}]$ पासून
$\Atom{\Obj P}{\Obj a}$ !!{derive} करू दिले पाहिजे. कदाचित हे निगमन नियम
म्हणूनही मांडायचे असेल. पण त्यासाठी ते शक्य तितक्या सर्वसाधारण शब्दांत
मांडता आले पाहिजे: फक्त $\Obj P$, $\Obj a$ आणि $\Obj v_0$ यांच्यासाठी
नव्हे, तर कोणत्याही !!{formula}~$!A$, पद~$t$ आणि !!{variable}~$x$
यांच्यासाठी. (!!{constant}s ही पदे असतात हे आठवा; पण !!{constant}s आणि
!!{function}s यांच्यापासून रचलेल्या अधिक जटिल पदांचाही आपण विचार करू.)
म्हणून आपल्याला असे काहीतरी म्हणता आले पाहिजे: ``जेव्हा तुम्ही
$\lforall[x][!A(x)]$ !!{derive} केले असेल, तेव्हा~$!A(t)$ चे अनुमान
करणे समर्थित असते---$\lforall[x]$ काढून आणि~$x$ च्या जागी~$t$ ठेवून
मिळणारे हे सूत्र आहे.'' पण ``$x$ च्या जागी~$t$ ठेवणे'' म्हणजे नेमके
काय? $!A(x)$ आणि~$!A(t)$ यांच्यातील संबंध कोणता? हे नेहमी चालते का?
%READERNOTE{OLFOL-004: गोठवलेल्या इंग्रजीतील पहिल्या सार्वत्रिक सूत्रात Atom चा दुसरा फलसाधक चुकीने सार्वत्रिक सूत्राबाहेर ठेवला आहे. पुढील सर्वसाधारण रूपाशी जुळण्यासाठी मराठीत P(v_0) हा पूर्ण आण्विक सूत्रभाग संख्यापकाच्या व्याप्तीत ठेवला आहे.}

हे काटेकोर करण्यासाठी आपण \emph{आदेशन} ही क्रिया परिभाषित करतो.
आदेशन प्रत्यक्षात गुंतागुंतीचे आहे, कारण~$!A$ मधील सर्व~$x$ फक्त~$t$ ने
बदलता येत नाहीत आणि प्रत्येक~$t$ चे कोणत्याही~$x$ साठी आदेशन करता येत
नाही. आपण हे पुन्हा विगमनाधारित व्याख्या वापरून हाताळू. पण हे
केल्यानंतर ``$\lforall[x][!A(x)]$ पासून $!A(t)$ चे अनुमान करा'' असा
निगमन नियम निर्दिष्ट करणे ही काटेकोर व्याख्या बनते. शिवाय,
$\lforall[x][!A(x)]$ पासून~$!A(t)$ तार्किकरीत्या निष्पन्न होते या अर्थाने
हा चांगला निगमन नियम आहे, हे दाखवता येईल. पण हे सिद्ध करण्यासाठी प्रथम
दर्शनी ``आपण हे का करत आहोत?'' असा प्रश्न पडेल अशी आणखी एक गोष्ट पुन्हा
सिद्ध करावी लागेल. $\lforall[x][!A(x)]$ पासून~$!A(t)$ तार्किकरीत्या
निष्पन्न होते हे पुढील तथ्यावर अवलंबून आहे: $\Sat{M}{!A(t)}$ लागू आहे
की नाही हे फक्त पद~$t$ च्या मूल्यावर अवलंबून असते. म्हणजे, $t$ ने
निर्देशित होणारा~$\Domain{M}$ मधील कोणताही !!{element} आपण $m$ मानला,
तर $\Sat{M}{!A(t)}[s]$ तेव्हा आणि केवळ तेव्हाच, जेव्हा
$\Sat{M}{!A(x)}[\Subst{s}{m}{x}]$. $t$ मध्ये !!{variable}s असली तरीही
हे लागू असते; पण हा परिणाम नेमका कसा मांडायचा याबाबत काळजी घ्यावी
लागेल.

\end{document}
